\section{Motivation and Context}
With more uncrewed aerial vehicles (UAVs) in the air each year, and a growing desire for them to operate in controlled airspaces, there is a concern regarding their ability to maintain separation from hazards in the air. United States and European federal agencies are developing policy that necessitates detect-and-avoid (DAA) systems for UAVs, which parallel the visual and instrument flight rules (VFR \& IFR) that dictate how pilots avoid traffic \cite{easa_easy_access_uas_2024} \cite{rtca_sc228_2025} \cite{serres_wu_2021_nasa_tm}.

While some proposed DAA systems utilise optical and infrared sensors \cite{serres_wu_2021_nasa_tm}, radar sensors may also be suitable, as they can perform sensing over great distances in all weather conditions, due to their relatively low atmospheric attenuation (compared with optical systems) \cite{borden_radar_imaging_airborne_targets_1999}. The inconvenient side effect of their long wavelength property is that angular resolution is much lower compared to similarly sized optical systems. If this challenge is overcome, the implementation of radar as the primary sensing mode or an additional sensing mode will improve the situational awareness, and thus the safety, of UAVs.

The main historical objection towards radar has not been with its usefulness, but with its fit within the SWaP constraints of smaller aircraft. This historical constraint is now weakening because both sides of the radar stack have miniaturised. Texas Instrument's AWR2243 is a single 76-81 GHz radar chip that integrates phase-locked loops (PLLs), three transmitters, four receivers, a mixer, an ADC, and self calibration, all on a 10.4mm \[\times\] 10.4mm printed circuit board (PCB)-friendly package. Higher frequency, lower noise, electronics and the miniaturisation of the transistor have been key enablers of this technology. Higher frequency radar allows of higher resolution imaging or smaller aperture; lower noise allows for an increase sensitivity to targets. On the compute side, modern UAV systems often have a dedicated, high compute companion computer in addition to the flight compute. NVIDIA's Jetson series computers such as Xavier NX, Orin Nano and Orin NX providing up 21, 67, 157 trillion operations per second (TOPS) on 8-bit integers respectively, while each weighting under 50g (without a heat sink). 

Plextek’s mmWave-based 60 GHz sense-and-avoid radar for UAVs demonstrates the feasibility of pairing miniaturised radar hardware with modern edge computing for the detection of people, vehicles, poles, and buildings at ranges of hundreds of metres. The wider UAV market shows a similar shift toward more computationally demanding onboard tasks. Skydio’s X10, for instance, incorporates both an NVIDIA Jetson Orin SoC and a Qualcomm QRB5165 SoC within a 2.11 kg platform while retaining up to 40 minutes of flight time, illustrating that substantial embedded compute is now compatible with practical UAV SWaP limits. These developments suggest that recent improvements in radar miniaturisation and embedded heterogeneous computing are promoting a new class of computation-heavy UAV sensing systems, making it timely to revisit radar target tracking and to determine the most suitable signal-processing pipeline for UAV detect-and-avoid applications. 

Despite the challenges, UAVs are the fastest growing market for radar systems, projected to grow with a compound annual growth rate (CAGR) of 11.22\% globally through 2030 \cite{mordor_airborne_radars_market_2025}. To make these emerging radar systems more performant, accessible, and safe, it would be valuable to find the most suitable radar signal processing pipeline for UAV DAA systems, revisiting the old problem of radar target tracking in the face of unique hardware conditions.

Software-defined architecture (SDAs) are increasingly preferred over hardware defined radar systems for UAV applications as software modifications can be shipped faster than hardware modifications \cite{mordor_airborne_radars_market_2025}. As a result, graphics processing units (GPUs) are favoured over non-programmable digital signal processors such as application specific integrated circuits (ASICs). GPU-accelerated implementations are also favourable for computationally intensive stages of the radar processing chain, as many core operations, such as fast Fourier transforms (FFTs) and matrix operations, are highly data-parallel, allowing for large increases in throughput when compared to serial processors. Additionally, GPU-based pipelines can leverage mature developer friendly, GPU software ecosystems and optimised libraries \cite{zaknoun2024_gpu_radar}. Additionally, the incorporation of AI into radar tracking algorithms has further encouraged these highly parallel, GPU-based pipelines \cite{electronics13214251}. Key open questions remain around balancing the compute saved through GPU parallelism against the cost of data movement and synchronisation, and how unified (shared) CPU--GPU memory on embedded platforms changes this trade-off. 

Given the miniaturisation of radar enabling technology, the potential benefits for radar base DAA systems, and the range limitations of systems suitable for small radar systems. This report proposes a software pipeline design that allows for the radar tracking of aerial hazards on UAVs with strict SWaP constraints. Particular emphasis will be placed on maximising the sensitivity to faraway targets, a requirement for operating in shared airspace, and minimising computation resources and latency. The knowledge gaps that will be addressed are detailed in the following section. 